\documentclass[a4paper,10pt]{article}
\usepackage[francais]{babel}
\usepackage{url}

%\input{letter-ensiie.sty}
\input{letter-icube-uds.sty}

\enCouleur

\begin{document}

\name{St\'ephane \textsc{Genaud}}
\fonction{Professeur ENSIIE, mise à disposition Université de Strasbourg}
\telephoneillkirch{+33 (0)3688 54542}
\email{genaud@unistra.fr}

\begin{letter}[Attestation et Recommendation de Luke \textsc{Bertot}]{à}%
{CNU 27, Campagne de Qualification 2020}

Chers et chères Collègues,

Dans le cadre de la candidature de  Luke Bertot à la qualification aux fonctions
de  maître   de  conférences,  j'apporte   par  cette  lettre   une  information
essentiellement qualitative sur l'activité de recherche de M. Bertot.  

Le domaine  de recherche  est celui des  infrastructures de  calcul distribuées,
plus  particulièrement le  cloud. J'ai  été le  directeur de  la thèse  que Luke
Bertot a commencé en octobre 2014 et  soutenue en juin 2019.  Cette thèse a donc
duré relativement longtemps. Sur le plan de l'enseignement, Luke Bertot a occupé
pendant ses deux dernières années de thèse un poste d'ATER 192h dans un IUT, qui
représente une  charge assez lourde  dans la phase  finale d'une thèse.   Sur le
plan  scientifique,  le  temps  passé   se  justifie  notamment  par  le  nombre
considérable   d'expérimentations   nécessaire    pour   obtenir   des   données
d'observation  du comportement  d'applications distribuées  sur des  plateformes
réelles. Nous avions  également décidé de confronter les  mesures observées avec
celles  prédites par  un  simulateur  utilisé par  la  communauté.  De  nombreux
ajustements ont  du être  apportés au simulateur  pour affiner  ses prédictions.
Enfin,  la  contribution   de  Luke  se  situe  dans   l'ajout  d'une  dimension
stochastique à la simulation de telles plateformes.

Au final, le travail  de thèse a sollicité de nombreuses  compétences de la part
de   Luke:   logicielles   pour   maitriser  les   applications   déployées   en
expérimentation et modifier  les simulateurs, système pour  programmer les plans
d'expériences sur  les plateformes,  et statistiques  pour construire  le modèle
stochastique. Je  considère que Luke s'est  parfaitement adapté à chacun  de ses
domaines. J'ai  pu apprécié sa curiosité  insatiable et sa persévérance  dans la
compréhension  des domaines  qu'il  a  du découvrir  pendant  sa  thèse. De  mon
expérience de  directeur de thèse,  Luke est un chercheur  scientifiquement très
solide. Dans toutes nos discussions ainsi que dans la rédaction d'articles, il a
toujours  avancé ses  idées de  manière  très argumentée,  démontrant ainsi  son
autonomie scientifique.

Pour  ces raisons,  je considère  que  Luke Bertot  a toutes  les qualités  pour
continuer une carrière dans le monde de l'enseignement et recherche.

Je vous prie, chers et chères collègues, d'agréer mes sentiments respectueux.
\end{letter}
\begin{flushright}
\includegraphics[width=.2\textwidth]{signgenaud.jpg}
\end{flushright}
\end{document}
